% 本文件由 LaTeX 排版 Agent 根据有效 translation.json 自动生成。
% author=Augustine of Hippo; work_id=1401; title=论天主教会之道德

\chapter{论天主教会之道德}

% structure: p0001 source=h1 target=omitted reason=page-title-already-rendered-by-parent

\Emph{撰写者：\Person{圣奥斯定}（\Place{希波}），公元388年。}\par

\Emph{开端即确立：教会圣洁生活的习俗应当归于人的至善，即天主。我们必须以最高的爱寻求天主；此教义在天主教会内由两约的权威所支持。四枢德从这种爱的不同形式得名。随后是关于爱近人的义务。在天主教会中，我们找到节德和真正基督徒生活的榜样。}\par

% structure: p0004 source=h2 target=section reason=relative-heading-rank; base=h2

\section{第一章 如何驳斥摩尼教徒的妄称。两项摩尼教的谬误}

1. 在我们的其他著作中，大概已经足够地回应了摩尼教徒对称为旧约的法律所作的无知和亵渎的攻击；他们以虚浮的夸耀和未受教者的赞同进行攻击。在此我也要简短地谈及这一主题。因为每个中等智力的人都能轻易看出，圣经的解释应当从那些公认的圣经教师那里寻求；而且可能发生，并且确实经常发生，许多事情对无知者显得荒谬，但被有学识者解释后，就显出其卓越，并且在解释中因着妨碍理解含义的障碍而更加令人喜悦地接受。就旧约的圣书而言，这种情况常发生，只要遇到困难的人求助于虔诚的教师，而非亵渎的批评者，并且他以寻求真理的渴望，而非鲁莽的反对开始探究。若求问者遇到一些主教、司铎或天主教会中任何官员或服务人员，他们或回避揭示奥秘，或满足于简单的信德而不渴求更深的学识，则求问者不应绝望于寻得真理的知识，因为并非所有被问者都能教导，也并非所有求问者都配得学习真理。勤勉和虔诚都是必要的：一方面，我们必须有知识以寻得真理；另一方面，我们必须配得获得知识。\par

2. 然而，由于摩尼教徒有两种捕捉粗心者的伎俩，以使他们视其为导师——其一是指责圣经（他们要么误解，要么希望被误解），其二是炫耀贞洁和显著的节制——本书将按公教教导包含我们关于生活和道德的教义，或许会使人看出，伪装德行何其容易，而拥有德行何其困难。我将尽力避免以他们攻击自己所不知之事的那种暴力，去攻击他们我所熟知的弱点；因为我希望他们若能痊愈，胜过被征服。我将引用他们必定相信的圣经见证，因为这些将取自新约；即使从新约中，我也将不取那些摩尼教徒在窘迫时惯称伪作的经文，而取他们不得不承认和认可的经文。对于每一宗徒教导的见证，我将从旧约中提出相似的陈述，以便若他们终愿从持续的梦幻中醒来，朝向基督信仰的光明上升，他们就能发现他们自称的生活何其远离基督信仰，而他们所挑剔的圣经何其真实地属于基督。\par

% structure: p0007 source=h2 target=section reason=relative-heading-rank; base=h2

\section{第二章 他依从摩尼教徒的错误方法，从论证开始}

3. 那么，我从何处开始？从权威，还是从推理？在自然的秩序中，我们学习任何事物时，权威先于推理。因为一种推理在给出之后，若需要权威来确认，就似乎软弱。但由于人的心智被黑暗的熟悉所蒙蔽，这黑暗在罪恶和恶习的夜晚遮盖他们，不能以适合理性清晰与纯净的方式感知，就有最有益的预备，在权威的合宜荫蔽下，将眩目的眼睛引入真理之光。但既然我们与那些在思想、言语和行为上全都乖谬的人打交道，且他们最坚持从论证开始，我便迁就他们，采取我认为在讨论中错误的方法。因为我喜欢尽可能模仿我主耶稣基督的温柔，祂亲自承担了死亡的邪恶本身，愿意从中解放我们。\par

% structure: p0009 source=h2 target=section reason=relative-heading-rank; base=h2

\section{第三章 幸福在于享有人的至善。至善的两个条件：第一，无物优于它；第二，不能违背意志而失去}

4. 那么，依照理性，人应当如何生活？我们当然全都渴望生活幸福；几乎在说出这句话之前，没有一个人不同意。但我认为，幸福这个称号既不属于没有拥有他所爱之物（无论是什么）的人，也不属于拥有他所爱之物若它是有害的人，也不属于不拥有他所爱之物若它有害的人，也不属于不爱他所拥有的、尽管那物是完美之善的人。因为寻求不能得到之物的人受折磨，得到不应得之物的人受欺骗，不寻求值得寻求之物的人有病态。在所有这些情形中，心智不能不痛苦，而幸福和痛苦不能同时存在于一人身上；所以，在这些情形中，无人能幸福。于是我找到第四种情形，其中幸福生活存在：当人的至善被爱且被拥有之时。因为我们所谓享乐，不正是拥有所爱之物吗？无人能在不享有人之至善的情况下幸福，也无人能享用它而不幸福。那么，如果我们想生活幸福，就必须拥有我们的至善。\par

5. 我们现在必须探求人的至善是什么；这当然不能是任何低于人本身的东西。因为凡是追求低于自身者，便变得低于自身。但每个人都应当追求那最好的。因此人的至善并不低于人。那么，它是否与人本身相似呢？倘若人所能享有的东西中没有高于人的，那必然如此。但如果我们发现某种东西既高于人，又能为爱它的人所拥有，那么谁还能怀疑，寻求幸福的人应当努力达到那比追求者本身更卓越的东西呢？因为如果幸福在于享受一种没有比它更好的善——我们称之为\OriginalTerm{至善}{summum bonum}——那么一个人若尚未达到他的至善，又如何能恰当地被称为幸福呢？或者，如果还有更好的东西等待我们去达到，那东西又怎能是至善呢？既然如此，至善必定是不能违背意志而失去的东西。因为没有人能对自己知道可能被夺走、尽管他愿意保持和珍视的善感到安心。但如果一个人对他所享有的善没有信心，他在害怕失去这种善的时候，又怎能幸福呢？\par

% structure: p0012 source=h2 target=section reason=relative-heading-rank; base=h2

\section{第四章 人——是什么？}

6. 那么，让我们来看看什么比人更好。这必然难以发现，除非我们首先询问并考察人是什么。现在，我并非被要求给人下一个定义。这里的问题在我看来是——既然几乎所有人都同意，或者至少足以说明的是，我现在与之打交道的人与我同感，认为我们由灵魂和身体组成——什么是人？他是两者兼备吗？或者仅仅是身体，或者仅仅是灵魂？因为虽然灵魂和身体是两样东西，而且没有另一方就不能被称为人（因为身体没有灵魂就不是人，同样，如果没有身体赋予其生命，灵魂也不是人），但仍然可能其中一样被认为就是人，并被称为人。那么，我们把什么称为人呢？是如同双驾马车或半人马一样的灵魂和身体吗？还是仅仅指身体，如同灵魂支配它、为它服务的那个身体，正如\Quote{灯}这个词所指的不是光与灯罩一起，而是仅指灯罩，然而它是因光而如此称呼的？还是仅指心灵，并且是因为它支配的身体而如此称呼，正如\Quote{骑手}一词不是指人与马，而是仅指人，并且是就他支配马而言的？这个争论不易解决；或者，如果证明是清楚的，陈述需要时间。这是时间和精力的消耗，我们无需承担。因为无论\Quote{人}这个名称属于两者还是仅属于灵魂，人的至善都不是身体的至善；而无论是灵魂和身体两者的至善，还是仅灵魂的至善，那才是人的至善。\par

% structure: p0014 source=h2 target=section reason=relative-heading-rank; base=h2

\section{第五章 人的至善并非只是身体的至善，而是灵魂的至善}

7. 现在如果我们问什么是身体的至善，理性迫使我们承认，那是使身体达到最佳状态的事物。但在所有使身体强健的事物中，没有什么比灵魂更好或更伟大的了。因此，身体的至善不是身体上的快乐，不是没有痛苦，不是力量，不是美丽，不是敏捷，或任何其他通常被算作身体之善的东西，而仅仅是灵魂。因为灵魂通过其临在赋予身体上述所有事物，而最重要的是生命。由此我得出结论：无论我们将人与身体一起称为灵魂和身体，还是仅称为灵魂，灵魂都不是人的至善。因为正如根据理性，身体的至善是比身体更好的东西，身体从中获得活力和生命；同样，无论灵魂本身是人，还是灵魂和身体两者，我们必须发现是否有任何先于灵魂本身的东西，灵魂追随它便能达到其自身种类所能达到的至善的完美。如果能够找到这样的东西，一切不确定性都必须终结，我们必须宣布这正是人的真正的至善。\par

8. 再者，如果身体就是人，那么必须承认灵魂是人的至善。但显然，当我们讨论道德时——当我们探问必须持守怎样的生活方式才能获得幸福时——诫命并非指向身体，我们讨论的也非身体的训练。简言之，良好\Emph{习俗}的遵守属于我们当中探问和学习的部分，而这是灵魂的特权；因此，当我们谈到获得\OriginalTerm{德性}{virtus}时，问题不涉及身体。但是，既然结果如下：被拥有德性的灵魂所统治的身体，统治得更好、更尊贵，并且由于公正地治理它的灵魂的完美而处于其最大的完美之中，那么使灵魂完美的东西便是人的至善，即使我们称身体为人。因为如果我的马车夫在听从我的情况下，以最令人满意的方式喂养和驾驶他所负责的马，而他自己因良好的行为而更多地享受我的赏赐，那么谁能否认，马的良好状况以及马车夫的良好状况都归功于我？因此，在我看来，问题不在于灵魂和身体是否是人，或仅仅是灵魂，或仅仅是身体，而在于什么使灵魂完美；因为一旦获得了这个，人就不能不成为完美的，或者至少比缺少这一样东西要好得多。\par

% structure: p0017 source=h2 target=section reason=relative-heading-rank; base=h2

\section{第六章 德性使灵魂达于圆满；灵魂随从天主而得德性；随从天主即是\OriginalTerm{幸福生活}{beata vita}}

9. 无人会质疑德行赋予灵魂完美。但一个非常合适的研究课题是：这德行能否独立存在，抑或只能存在于灵魂中？这里再次出现一个深奥的讨论，需要长篇论述；但或许我的概要能达其目的。我信靠天主的助佑，使我们尽管软弱，却能简明而清晰地就这些重大课题施教。无论哪种情况——德行能否脱离灵魂而独立存在，抑或只能存在于灵魂中——无疑，在追求德行时，灵魂是在追随某物，而此物要么是灵魂自身，要么是德行，要么是其他东西。但如果灵魂在追求德行时追随自身，它就是在追随一件愚蠢之物；因为获得德行之前，它是愚蠢的。追随者的欲望之巅，正是达到他所追随的目标。因此，灵魂要么不愿达到它所追随的目标——这完全荒谬无理——要么，在追随自身而仍愚蠢时，它达到了它所逃避的愚蠢。但如果它追求德行以渴望达到它，它怎能追随不存在之物？或者，它怎能渴望达到它已拥有之物？因此，要么德行存在于灵魂之外，要么，如果我们只允许将德行之名赋予那只能存在于灵魂中的、智慧灵魂的习惯与性情，那么我们必须承认，灵魂是为了在自身中产生德行而追随其他某物；因为据我推理，灵魂既不能因追随虚无，也不能因追随愚蠢而获得智慧。\par

10. 那么，这其他某物——灵魂因追随它而获得德行与智慧——要么是智慧人，要么是天主。但我们已经说过，它必须是我们不能违背自己意愿而失去的。无人会认为有必要询问：一位智慧人，假设我们满足于追随他，是否可能在我们不情愿或执着的情况下被夺走？那么，留下的就是天主；我们追随天主而生活良善，达到天主而生活既良善又幸福。若有人否认天主的存在，我何必考虑与他们辩论的方式，因为甚至不确定是否应当与他们辩论？无论如何，这需要不同的出发点、不同的计划、不同于我们目前所从事的研究。我现在是对那些不否认天主存在、并承认天主并不忽视人类事务的人说话。因为我想，凡有任何宗教信仰的人，都会主张神圣眷顾至少照料我们的灵魂。\par

% structure: p0020 source=h2 target=section reason=relative-heading-rank; base=h2

\section{第七章——从圣经获得对天主的认识。神圣救赎计划的纲要与主要奥迹}

11. 但我们怎能追随我们看不见的那位？我们，不仅是人，而且是理解力软弱的人，又怎能看见祂？因为虽然天主不是以肉眼，而是以心灵被看见，但哪里能找到这样的心灵，它在被愚蠢遮蔽时，竟能成功，甚至尝试，去吸收那光明？因此，我们必须求助于那些我们有理由认为是智慧之人的教导。至此，推理引导我们至此。因为在人事中，推理被使用，并非因其更确定，而是因其更易于运用。但当我们来到神圣事物前，这能力便转向；它不能凝视；它喘息，气促，燃起渴望；它从真理之光中退回，再次转回它惯常的黑暗，并非出于选择，而是由于疲惫。这是何等可怕的灾祸：灵魂在寻求从劳苦中安息时，反而陷入更大的无助！因此，当我们急忙隐退到黑暗中时，最好是由可钦崇的智慧所安排，让我们遇见权威的友善荫庇，并被其内容的奇妙特征以及其书卷的言辞所吸引——这些言辞如同影子，象征并调适真理。\par

12. 为我们的救恩还能做什么更多的事？有什么比神圣眷顾更恩慈、更慷慨的呢？它当人背离其律法，并因贪恋必死之物而受到公正报应、生出必死后裔时，仍没有完全抛弃他。因为在这最公义的治理中——其途径奇异而不可测——借着受造物中未知的连结，既有惩罚的严厉，也有救恩的仁慈。这一切何其美丽，何其伟大，何其堪配天主，总之，何其真实——这正是我们所寻求的一切——我们永远无法领会，除非从眼前的人事开始，持守真宗教的信德与诫命，继续行走在天主为我们所预备的道路上，不偏左右。这道路由圣祖的分别、法律的纽带、先知的预见、宗徒的见证、殉道者的鲜血以及外邦人的归顺所保障。因此，从这一点起，请勿有人向我询问我的意见，而是让我们聆听神谕，并将我们软弱的推论屈从于天主的宣告。\par

% structure: p0023 source=h2 target=section reason=relative-heading-rank; base=h2

\section{第八章——天主是至善，我们当以最高挚爱寻求祂}

13. 让我们看看主自己在福音中如何教导我们生活；也同样看看宗徒保禄如何教导——因为摩尼教徒不敢拒绝这些圣经。让我们听，哦基督，您指示给我们最终目标是什么；那显然是我们被告知要以最高情感追求的目标。祂说：\Quote{你当爱上主你的天主}。也请告诉我，我祈求您，爱的尺度是什么；因为我担心我心中燃起的渴望在热忱上会过度或不足。祂说：\Quote{以你的全心}。这还不够。\Quote{以你的全灵}。这仍不够。\Quote{以你的全意}。\BibleRef{玛 22:37} 你还有什么想要的？也许，如果我能看到更多的可能，我会想要更多。保禄对此怎么说？他说：\Quote{我们知道，凡爱天主的人，一切都为他们效力}。也让他说出爱的尺度。他说：\Quote{那么，谁能使我们与基督的爱隔绝呢？是困苦吗？是窘迫吗？是迫害吗？是饥饿吗？是赤身吗？是危险吗？是刀剑吗？} 因此，我们听到了要爱什么以及要爱多少；我们必须追求这点，并将我们所有的计划都指向这点。我们一切美善的完美和完美的美善就是天主。我们既不能不足也不能过度：前者危险，后者不可能。\par

% structure: p0025 source=h2 target=section reason=relative-heading-rank; base=h2

\section{第九章 论新旧约在爱德诫命上的一致}

14. 现在，让我们检查，或者更确切地说，让我们留意——因为这是显而易见的，可以立刻看到——旧约的权威是否也与那些取自福音和宗徒的陈述一致。有什么必要谈论第一个陈述呢？因为所有人都清楚它是引自梅瑟所颁布的法律。因为那里写道：\Quote{你当爱上主你的天主，以你的全心，以你的全灵，以你的全意}。\BibleRef{申 6:5} 而且不必再往前去找旧约中与宗徒相比较的一段，他自己已经加上了一段。因为他说过，没有任何困苦、窘迫、迫害、身体匮乏的压力、危险、刀剑能使我们与基督的爱隔绝之后，他立刻加上：\Quote{正如经上记载：\InnerQuote{为了你，我们终日受磨难；我们被视为待宰的羊。}} 摩尼教徒习惯说这是篡改——他们如此无法回应，以至于在绝境中被迫这样说。但每个人都可以看出，当人们被证明是错误的时，这是他们唯一能说的话。\par

15. 然而我问他们，他们是否否认这在旧约中说过，或者他们是否认为旧约中的那段与宗徒的话不一致？首先，书籍会证明它；其次，那些偏离正轨的诡辩者，如果他们稍加反思并考虑所说的话，就会同意我的看法；否则，我将迫使他们接受公正判断者的意见。还有什么能比这些段落更和谐一致呢？因为在此生中，困苦、窘迫、迫害、饥饿、赤身、危险给人带来极大的痛苦。所以所有这些话都包含在来自法律的单一引文中，那里说：\Quote{为了你，我们受苦}。唯一剩下的是刀剑，它并不施加痛苦的生命，而是夺去它所遇到的任何生命。与此相应的是：\Quote{我们被视为待宰的羊}。爱不能再比\Quote{为了你}这话表达得更清楚了。那么，假设这个证言不在宗徒保禄那里，而是由我引用，你这个异端者，难道你不必须证明要么这在旧约中没有记载，要么它与宗徒不一致吗？如果你不敢说这两点中的任何一点（因为你被手稿的朗读所困，这会显示它是写着的，又被常识所困，它看出没有比宗徒所说的更一致的了），你为什么还认为指控圣经被篡改有什么力量呢？再者，你会对那个对你说以下话的人作何回答：\Quote{这就是我所理解的，这是我的观点，这是我的信仰，我阅读这些书仅仅是因为我看到书中的一切都与基督信仰一致？} 或者立刻告诉我，你是否会故意当面告诉我，我们不该相信宗徒和殉道者被认为是为主基督忍受了极大的苦难，并被他们的迫害者视为待宰的羊？如果你不能说这个，你为什么要控告那本书——在书中我找到了你承认我应该相信的东西？\par

% structure: p0028 source=h2 target=section reason=relative-heading-rank; base=h2

\section{第十章 论教会关于天主的教导。摩尼教徒的两个神}

16. 你会说，你承认我们必须爱天主，但不是崇拜那些承认旧约权威的人所崇拜的天主吗？那样的话，你拒绝崇拜那创造天地的天主，因为那正是这些书中一直展现的天主。而你承认整个宇宙（称为天和地）有一位天主，且是一位善的天主作为其创造者和建立者。因为在对你谈论天主时，我们必须区分。因为你认为有两个神，一个善，一个恶。\par

但若你们说，你们崇拜并赞同崇拜那创造天地的天主，却不崇拜那受\WorkTitle{旧约}权威所支持的天主，那你们就是无礼地（尽管是徒劳地）试图将与我们真正持有的有益而健全的教义全然不同的观点和意见归咎于我们。你们愚昧而亵渎的言谈，也根本无法与那些博学而虔诚的天主教会人士向愿意且配得的人解释这些圣经的注释相比。我们对法律和先知的理解与你们所猜想的大相径庭。请不要再误会我们。我们不崇拜一位后悔、嫉妒、有所缺乏、残忍、以人或兽的血为乐、以罪恶和过犯为乐、或其对大地的占有仅限于一小角落的天主。这些以及诸如此类的愚蠢观念，正是你们惯于长篇大论地谴责的。你们的谴责并不触及我们。你们以可笑之极的激烈态度攻击老妇或孩童的幻想。凡被你们以此方式说服而加入你们的人，所表明的并非教会教义有什么缺陷，而只是他自身对教义的无知。\par

17. 因此，倘若你们尚有一丝人情——倘若你们还关心自己的福祉——你们就该以勤奋和虔诚来审视这些\WorkTitle{圣经}章节的含意。你们这些不幸的人啊，应当审视；因为我们同样严厉而充分地谴责任何将不配于天主之事归诸天主的信仰，并且对于那些按字面理解这些段落的人，我们纠正其无知的错误，而若坚持此错误，我们则视之为荒谬。此外，在你们无法理解的许多其他事上，天主教的教导对那已超越心智孩提阶段（并非按年龄，而是按知识和悟性——在迈向智慧上已年迈）之人的信仰有一种约束。因为我们得知，相信天主受任何空间（即使是无限的空间）所限是愚蠢的；并且认为天主或祂的任何部分从一个地方移动到另一个地方，是被视为不合法的。若有人设想天主的本体或本性中任何东西会遭受变化或转变，他将被视为犯有极度的亵渎。因此，我们中间有孩童般的人认为天主具有人形，他们以为天主实际上具有此形，这是一种极其卑劣的想法；也有许多成年人的心智中，天主的尊威以其不可侵犯和不变性显现，不仅超越人的身体，而且超越他们自己的心智本身。如我们所说，这些年龄的区分不在于时间，而在于德行和明辨。而在你们中间，没有人会将天主描绘成人形；但也没有人能使天主免于人类错误的玷污。至于那些像啼哭的婴儿一样在天主教会的怀抱中被喂养的人，如果他们没有被异端掳走，就会按照各自的力量和容量得到滋养，并最终以不同的方式达到完美的成人，进而达到智慧的成熟和苍苍白发，那时他们就可如其所愿地获得生命，获得圆满幸福的生命。\par

% structure: p0032 source=h2 target=section reason=relative-heading-rank; base=h2

\section{第十一章——天主是爱的唯一对象；因此祂是人的至善。无物优于天主。天主不会违背我们的意愿而失去}

18. 追随天主就是对幸福的渴望；达到天主就是幸福本身。我们因爱天主而追随祂；我们达到祂，并非成为祂之所是，而是通过接近祂，通过奇妙而无形地接触祂，并通过内在被祂的真理和圣善所光照和占有。祂是光本身；我们从祂那里得到光照。因此，导向幸福生活的最大诫命，也是第一条诫命就是：\BibleQuote{你当全心、全灵、全意爱上主你的天主。}因为爱主的人，万事都为他们得益。因此\Person{保禄}稍后补充说：\BibleQuote{我深信：无论是死亡，是生活，是天使，是掌权者，是现存的，是将来的，是高天，是深渊，或是任何别的受造物，都不能使我们与天主的爱相隔绝，这爱是在我们的主基督耶稣内的。}\BibleRef{罗 8:38} 那么，如果爱天主的人万事都得益，并且，正如无人怀疑的那样，至高的或完美的善不仅要被爱，而且要爱到没有任何事物比它更被爱，正如\BibleQuote{全心、全灵、全意}这句话所表达的，我请问，当这些事情都已确定且最坚定地相信时，谁不会立刻得出结论：我们优于一切事物而必须尽快达到的至善，不是别的，正是天主？而且，如果没有任何事物能使我们与祂的爱相隔绝，那么这岂不是比其他任何善更确实、更美好吗？\par

19. 但我们分开来考虑各点。没有人能以死亡威胁将我们从这爱中分离。因为我们用来爱天主的那个东西，除非不爱天主，否则不能死亡；死亡就是不爱天主，也就是我们在情感和追求中偏爱任何事物胜过祂。没有人以应许生命将我们从这爱中分离；因为没有人能用应许水将我们从泉源分离。天使不能分离我们；因为依附于天主的灵魂在力量上并不亚于天使。\OriginalTerm{德能}{virtue}不能分离我们；因为若这里所谓的德能是在这世界有权势的东西，那么依附于天主的灵魂远在全世界之上。或者，若这德能是我们灵魂本身的完美正直，那么在他人身上，它将有利于我们与天主的合一，而在我们自己身上，它本身将使我们与天主合一。目前的患难不能分离我们；因为我们越紧紧抓住那试图将我们与祂分离的那一位，我们就越不感觉到患难的重担。对将来之事的应许不能分离我们；因为各种将来的福乐在天主的应许中是最确定的，而且没有比天主本身更好的了，祂无疑已经临于那些真正依附于祂的人。高深之事不能分离我们；因为若所谓的高深之事是指知识的高深，我宁愿不追根究底，也不愿与天主分离；任何能消除错误的教导也不能使我与祂分离，因为与祂分离正是使人陷入错误的原因。或者，若所谓的高深之事是指这世界的上层和下层，那么对天堂的应许怎能使我与创造天堂的那一位分离呢？或者，从下面来的谁能恐吓我离弃天主呢？因为若非离弃祂，我本不会知道下面的事。总之，有什么地方能将我从祂的爱中移开呢？因为祂若不被任何地方所包含，就不能处处都在。\par

% structure: p0035 source=h2 target=section reason=relative-heading-rank; base=h2

\section{第十二章——我们借着爱，在服从天主中与祂合一}

20. \Quote{没有别的受造物，}他说，\Quote{能将我们分离。}深奥奥秘的人啊！他认为说\Quote{没有受造物}还不够，而是说\Quote{没有别的受造物}，以此教导我们：我们用来爱天主并借以依附天主的那个东西，即我们的灵魂和理智，本身就是一个受造物。因此，身体是另一个受造物；如果灵魂是理智所能感知的，并仅由此而被认识，那么别的受造物就是所有感官的对象，它们仿佛通过眼睛、耳朵、嗅觉、味觉或触觉来使自己被认识，而这必定低于仅凭理智所感知的东西。既然天主也可以被那配得的人凭借理智来认识，尽管祂远高于理智的受造物，因为祂是理智的创造者和作者，就存在一种危险：人的灵魂，由于被归入无形和非物质的事物之列，可能被认为与创造它的那一位具有\Emph{相同的}本性，因而因骄傲而从祂那里堕落，而它本应借着爱与祂合一。因为灵魂变得像天主，其程度取决于它服从天主以获得启发和光照的程度。如果它因那产生相似的服从而获得最大的亲近，那么它必定因那试图使相似性更大的僭越而远离祂。正是这种僭越，导致灵魂拒绝服从天主的法律，渴望像天主一样自主。\par

21. 因此，灵魂离天主越远——不是在空间上，而是在情感和对低于祂的事物的贪欲上——就越充满愚昧和可怜。所以，它借着爱回归天主——这种爱不是将灵魂置于与天主同等，而是置于祂之下。在这爱中，越热烈和急切，灵魂就越幸福和崇高，并且以天主为唯一的统治者，它将享有完全的自由。因此，它必须知道它是一个受造物。它必须相信这个真理：它的创造者永远拥有不可侵犯和不变的真与智慧的本性，并且，即使面对它渴望从中解脱的错误，它也必须承认它容易陷入愚昧和虚假。然而，它还必须小心，不要因对别的受造物（即这个可见的世界）的爱，而与那圣化它以达到永恒幸福的天主本身的爱分离。所以，没有别的受造物——因为我们自己就是一个受造物——能将我们从那在基督耶稣我们的主内的天主之爱中分离。\par

% structure: p0038 source=h2 target=section reason=relative-heading-rank; base=h2

\section{第十三章——我们借着基督和祂的圣神与天主不可分离地合一}

22. 让这位保禄告诉我们，谁是这位基督耶稣我们的主。\Quote{我们给那蒙召的，}他说，\Quote{宣讲基督是天主的德能，天主的智慧。}\BibleRef{格前 1:23}基督自己岂不是说：\BibleQuote{我是真理}\BibleRef{若 14:6}吗？因此，如果我们问什么是善生——即借着善生追求幸福——那必定是爱德能、爱智慧、爱真理，并全心、全灵、全意地爱那不可侵犯、永不改变的德能，那从不让位于愚昧的智慧，那没有变化或改变其恒定特性的真理。借着这德能、智慧和真理，圣父本身就被看见；因为经上说：\BibleQuote{除非经过我，没有人能到父那里去。}我们借着圣化而依附于祂。因为当我们被圣化时，我们燃起圆满和完美的爱，这是确保我们不离开天主、并且使我们效法祂而不是效法这世界的唯一保障；因为同一位宗徒说：\BibleQuote{祂预先定了我们，使我们效法祂圣子的肖像。}\BibleRef{罗 8:29}\par

23. 于是，我们藉着爱而肖似天主；藉此肖似、形构和从这世界的割离，我们不被那些本当属于我们的事物所混淆。这是由圣神成就的。他说：\BibleQuote{望德不会使我们蒙羞，因为天主的爱，藉着所赐与我们的圣神，已倾注在我们心中了。}\BibleRef{罗 5:5}但除非圣神本身常是完美不变的，我们不可能藉祂恢复完美。这显然只能发生在祂具有天主的性体和同一实体时，因为唯独天主常保不变与恒常。并非由我，而是由保禄断言：\BibleQuote{受造物被屈服于虚妄之下。}\BibleRef{罗 8:20}凡屈服于虚妄的，不能使我们脱离虚妄而结合于真理。但圣神为我们成就了这事。因此祂不是受造物。因为凡存在的，要么是天主，要么是受造物。\par

% structure: p0041 source=h2 target=section reason=relative-heading-rank; base=h2

\section{第十四章　我们藉爱依附天主圣三，我们的至善}

24. 因此我们应当爱天主，即天主圣三，圣父、圣子、圣神；因为必须说这就是天主本身，因为论及天主，真实而最崇高地说：\BibleQuote{万有出于祂，万有藉祂，万有在祂内。}这是保禄的话。他又加上什么？\BibleQuote{愿光荣归于祂。}\BibleRef{罗 11:36}这一切完全真实。他没有说\Quote{归于他们}；因为天主是唯一。所谓\Quote{愿光荣归于祂}，不就是愿至高、圆满、广泛的赞美归于祂吗？因为赞美越提升、越扩展，爱和情感就越热烈。当此情形，人类必然坚定稳妥地迈向完美和幸福的生命。我想，这就是当我们论及人的至善时所要寻求的一切，在生活和行为中一切都应归向它。因为善是明明存在的；我们已尽力通过推理，并藉超越我们推理的神圣权威表明，这善无非是天主本身。因为除了依附祂而获得幸福之外，还有什么能成为人的至善呢？这至善就是天主，我们只能藉情感、渴望和爱来依附祂。\par

% structure: p0043 source=h2 target=section reason=relative-heading-rank; base=h2

\section{第十五章　四枢德之基督徒定义}

25. 关于导引我们至幸福生活的德性，我认为德性无非就是对天主的完美之爱。我以为四枢德之分乃取自爱的四种形式。对于这四种德性（但愿众人在心中感受到它们的影响，如同在口中有其名一般！），我会毫不犹豫地如此定义：节德是爱将它自身完全献给所爱者；刚毅是爱为所爱者欣然承担一切；义德是爱只服务于所爱者，因而正确统治；智德是爱以明辨分辨何者阻碍它、何者帮助它。这爱的对象不是任何事物，而唯独是天主，至善、至高的智慧、完美的和谐。因此我们可以这样表达定义：节德是爱为天主保持自身完整无瑕；刚毅是爱为天主欣然承担一切；义德是爱只侍奉天主，因此妥善治理其他一切，如同它们隶属于人；智德是爱正确分辨什么帮助它归向天主，什么可能阻碍它。\par

% structure: p0045 source=h2 target=section reason=relative-heading-rank; base=h2

\section{第十六章　旧约与新约的和谐}

26. 在我依次简述按这些德性的生活方式之后，正如我所应许的，我将提出与我所引新约经文平行的旧约经文。因为难道只有保禄说我们应结合于天主，以致中间没有任何隔阂吗？先知不是最恰当而简洁地用这话说出同样的事吗？\Quote{我亲近天主，于我有益。}难道这一个词\Emph{亲近}不就表达了宗徒关于爱所长篇论述的一切吗？而\Quote{于我有益}不正是指向宗徒的话：\Quote{天主使一切协助那些爱祂的人获得益处}？如此，先知用一句话、两个词表明了爱的能力和果实。\par

27. 宗徒说天主子是天主的德能和天主的智慧——德能被理解为关涉行动，智慧关涉教导（正如福音中这两件事由下列话语表达：\Quote{万物是藉着祂而造的}，这属于行动和德能；然后论及教导和认识真理，祂说：\BibleQuote{生命是人的光}\BibleRef{若 1:3}）——有什么比旧约论及智慧的话更与这些经文一致呢？\Quote{智慧以德能达至地极，以温柔治理万有。}因为\Quote{以德能达至}表达德能，而\Quote{以温柔治理}表达技巧和方法。但如果这似乎隐晦，请看下文：\Quote{天主爱她，因为}他说，\Quote{她传授天主的知识，拣选祂的工程。}这里关于行动不再多提；因为\Quote{拣选工程}并非\Quote{工作}，所以这属于教导。行动仍有待与德能相应，以完成我们欲证实的真理。那么请读接下去的话：\Quote{但如果}他说，\Quote{生活中渴望的占有是可贵的，有什么比运作万有的智慧更可贵呢？}有什么比这更引人注目、更清晰，甚至更充分表达的呢？或者，如果你想要更多，请听另一段意义相同的话：\Quote{智慧}他说，\Quote{教导节制、正义和德能。}我想，\Quote{节制}指认识真理或教导；\Quote{正义和德能}指工作和行动。而我认为，在行动中的效能和默观中的节制——即天主的德能和天主的智慧，亦即天主子，赐给爱祂的人——这两者无与伦比。当同一位先知继续显明其价值时，就是这样陈述的：\Quote{智慧教导节制、正义和德能，在人的生活中没有比这些更有用的了。}\par

28. 也许有人会以为这些段落并非指天主圣子。那么，下文教导的是什么呢？\BibleQuote{她显示了自己出身的高贵，与天主同住}？\BibleRef{智 8:3} 出身若非指血统，又指什么？与圣父同住岂不要求并宣告平等？再者，正如保禄所说，天主圣子就是天主的智慧，\BibleRef{格前 1:24} 而主自己说：\BibleQuote{除了独生子，没有人认识父}，\BibleRef{玛 11:27} 还有什么比先知的话更一致的呢？\BibleQuote{智慧与祢同在，她认识祢的化工，在祢创造世界时已在场，也知道什么中悦祢的眼目。}\BibleRef{智 9:9} 又如基督被称为真理，这也可以从祂被称为父的光辉\BibleRef{希 1:3}（因为太阳周围除了从它发出的光辉，别无他物）得到证实，那么旧约中有什么比\BibleQuote{祢的真理环绕祢}这句话更清楚明白地与这相符呢？再者，智慧自己在福音中说：\BibleQuote{除非藉着我，没有人能到父那里去}；\BibleRef{若 14:6} 先知说：\BibleQuote{除非祢赐予智慧，谁能知道祢的旨意？}稍后又说：\BibleQuote{人们学习了中悦祢的事，并藉智慧得到了医治。}\BibleRef{智 9:17}\par

29. 保禄说：\BibleQuote{天主的爱藉着所赐予我们的圣神，已倾注在我们心中。}\BibleRef{罗 5:5} 先知说：\BibleQuote{知识的圣神必远避诡诈。}\BibleRef{智 1:5} 因为哪里有诡诈，哪里就没有爱。保禄说我们\BibleQuote{被模成天主圣子的肖像}，\BibleRef{罗 8:29} 先知说：\BibleQuote{祢面容的光辉印在我们身上。}保禄教导圣神就是天主，因此不是受造物；先知说：\BibleQuote{祢从高处差遣祢的圣神。}\BibleRef{智 9:17} 因为唯有天主是至高者，无有高于祂者。保禄表明天主圣三是唯一的天主，他说：\BibleQuote{愿荣耀归于祂}；\BibleRef{罗 11:36} 旧约中则说：\BibleQuote{以色列！你要听，上主你的天主是唯一的天主。}\BibleRef{申 6:4}\par

% structure: p0050 source=h2 target=section reason=relative-heading-rank; base=h2

\section{第十七章  劝勉摩尼教徒，呼吁他们悔改}

30. 你们还想要什么？为何愚昧固执地抗拒？为何用你们有害的教导败坏未经训练的心灵？两约的天主是同一位。因为从两个约中所引用的段落如此一致，其余部分亦然——只要你们愿意仔细公正地思考它们。但因许多表达方式卑下，适切于匍匐在地上的心灵，使他们藉属人的事升至属神的事；而许多表达方式是比喻性的，使探索的心从努力寻找其意义中获得更多益处，并在找到时获得更大喜悦；你们却歪曲圣神这令人赞叹的安排，为了欺骗和诱捕你们的追随者。至于天主眷顾为何容许你们这样做，以及宗徒所说的真理：\BibleQuote{必须有许多异端，好使那些被认可的人得以在你们中间显明出来}，\BibleRef{格前 11:19} 讨论这些事需要很长时间，而你们——我们目前所面对的人——无法理解。我深知你们。你们以物质形象为食，心灵粗俗病弱，却来思考远远高于你们所设想的神圣事物。\par

31. 因此，对于你们，我们必须努力的不是使你们理解神圣事物——这是不可能的——而是使你们渴望理解。这出于纯洁无诡诈的对天主的爱，这爱主要显在行为中，我们已经谈了许多。这爱由圣神所激发，引领到圣子，即天主的智慧，藉着祂，圣父本身得以被认识。因为若不全心竭力追求智慧和真理，就不可能找到。但当人按当得的方式寻求时，它不能从爱慕者那里隐退或隐藏。因此它的这些话——你们也惯常重复的——：\BibleQuote{你们求，就会得到；找，就会找到；敲门，就会给你们开。}\BibleRef{玛 7:7} \BibleQuote{没有隐藏的事不被显露出来的。}\BibleRef{玛 10:26} 是爱求，爱找，爱敲门，爱显露，爱又使人在所显露之事中持续。这种对智慧的爱和勤勉的探求，旧约并不禁止我们——如你们常假称的那样——反而极力敦促我们去行。\par

32. 那么，请你们细听并思考先知所说的话：\BibleQuote{智慧是光荣的，永不衰败；凡爱慕她的，她很容易见到；凡寻求她的，她必找到。她预先显示给渴望她的人，先使他们认识她。清晨寻求她的人，不必辛劳，因为必发现她坐在门前。因此，想念她即是智慧的圆满；凡警醒为她的人，必很快无忧无虑。因为她巡行各处，寻找配得上她的人，在道路上亲切地显示自己，在每次思念中迎接他们。因为智慧的真开端是渴望受教；受教之心即是爱；爱即是遵守她的法律；留心她的法律即是永不朽坏的保证；永不朽坏便使我们亲近天主。所以，对智慧的渴望引人进入王国。}\BibleRef{智 6:12} 你们还要顽固地敌视这些事吗？这些虽尚未理解，但难道不是使每个人都很清楚它们包含深奥不可言宣的道理吗？但愿你们能理解此处所说的话！你们就会立刻弃绝所有愚蠢的传说和空洞的物质想象，以极大的热情、真诚的爱和全备的确信，投身于至圣的慈母教会之怀。\par

% structure: p0054 source=h2 target=section reason=relative-heading-rank; base=h2

\section{第十八章  惟在天主教会中，两约和谐，完美真理得以建立}

33. 照我微薄的能力，我本可逐点剖析，阐述并证明我所学得的真理——这些真理通常比言语所能表达的更为卓越崇高；但当你们对着它狂吠时，这事便无法进行。因为经上并非徒然说：\BibleQuote{不要把圣物给狗。}\BibleRef{玛 7:6} 不要发怒。我也曾狂吠过，也曾是条狗；那时，正如所应得的，我得到的不是教导的食粮，而是管教的棍棒。但若你们内有我们此刻所谈论的爱，或有一天，当你们有了与所当认识的真理之伟大相称的爱时，愿天主俯允指示你们：在摩尼教徒中，既没有那导向智慧与真理之巅的基督徒信仰——达到这信仰即是真正的幸福生命——这信仰也不在任何其他地方，唯独存在于天主教的教导中。宗徒保禄岂不是在渴望这一点吗？他说：\BibleQuote{因此，我在圣父面前屈膝——天上地下的一切家族都由祂得名——愿祂照着祂荣耀的丰富，藉着祂的圣神，以大能坚固你们内在的人，使基督因着信德住在你们心中，叫你们在爱德上根深蒂固，得以和众圣徒一同领悟何为高、深、长、阔，并认识基督那超越知识的爱，使你们充满天主的一切圆满。}\BibleRef{弗 3:14} 还有什么能比这表达得更清楚呢？\par

34. 请你们稍微清醒，我恳求你们，看看两经典籍的和谐，它清楚明确地指明我们行事为人的方式应当如何，以及万事万物应归向何处。福音劝勉我们爱天主，当它说：\BibleQuote{求，寻找，敲门}\BibleRef{玛 7:7}；保禄也劝勉我们，当他说：\BibleQuote{叫你们在爱德上根深蒂固，得以领悟}\BibleRef{弗 3:7}；先知同样说，智慧易于为那些爱慕它、寻求它、渴望它、警醒它、思念它、关心它的人所认识。两经典籍的和谐指出了心灵的得救和幸福之路；而你们却宁愿对这些事实狂吠，也不愿顺从它们。我将用一句话告诉你们我的想法。请你们以同样平和的性情和同样的热忱，去聆听天主教会中的学者，就像我曾花了九年时间跟随你们时那样；那么，就不需要那么长的时间——在你们愚弄我的那段时间里——你们就会看到真理与虚妄之间的区别，而且时间会大大缩短。\par

% structure: p0057 source=h2 target=section reason=relative-heading-rank; base=h2

\section{第十九章——论节德的职责，按圣经的论述}

35. 现在该回到四种德行，并逐一根据它们来拟定并规定一种生活方式。首先，让我们思考节德，它应许我们在那使我们与天主结合的爱中，获得某种完整与不朽。节德的职责在于约束和平息那些使我们渴望远离天主法律、远离对祂美善的享受——简言之，远离幸福生命的事物。因为真理的居所就在那里；在享受对真理的默观并紧紧依恋它时，我们确实是幸福的；但背离它，人们就会陷入巨大的错误和痛苦。因为宗徒说：\BibleQuote{贪欲是万恶的根源；有些人因追随贪欲，在信德上遭了覆舟之祸，并使自己受了许多刺心的痛苦。}\BibleRef{弟前 6:10} 对正确理解的人来说，灵魂的这种罪在旧约中藉着亚当在乐园中的悖逆已被十分清楚地展示。因此，正如宗徒所说：\BibleQuote{在亚当内众人都死了，同样，在基督内众人都要复活。}\BibleRef{格前 15:22} 哦，这些奥迹何其深奥！但我克制自己，因为我现在不是要教导你们真理，而是要纠正你们的错误，如果我能的话——也就是说，如果天主为你们而助成我的目的。\par

36. 保禄于是说贪欲是万恶的根源；旧约也藉此暗示第一个人因贪欲而堕落。保禄要我们\BibleQuote{脱去旧人，穿上新人}\BibleRef{哥 3:9}。所谓旧人，他指那犯了罪的亚当；所谓新人，他指天主之子为救赎我们而圣化时所取的那个人性。因为他在另一处说：\BibleQuote{第一个人出于地，属于土；第二个人出于天，属于天。那属于土的怎样，凡属于土的也怎样；那属于天的怎样，凡属于天的也怎样。我们既带着那属于土的形象，也要带着那属于天的形象}\BibleRef{格前 15:47}——也就是说，脱去旧人，穿上新人。于是，节德的全部职责便是脱去旧人，并在天主内更新——也就是说，藐视一切肉体的快乐和世俗的喝采，将全部的爱转向神圣而不可见的事物。因此便有了那句极令人赞叹的经文：\BibleQuote{我们外在的人虽在朽坏，内在的人却日日更新。}\BibleRef{格后 4:16} 请听先知歌唱：\BibleQuote{天主，求你给我再造一颗纯洁的心，在我内里更新一个坚定的精神。}面对这样的和谐，除了瞎眼的狂吠者，还能说什么呢？\par

% structure: p0060 source=h2 target=section reason=relative-heading-rank; base=h2

\section{第二十章——我们应鄙视一切可感之物，唯独爱慕天主}

37. 肉身的快乐源于一切与肉身感官接触的事物，这些事物有人称之为感觉对象；其中最高贵的是光，按这个词的通常含义，因为在我们感官中，灵魂藉着肉身行动时所用的，没有比眼睛更宝贵的，因此在圣经中，一切感觉对象都被称为可见之物。新约中，我们被以下话语警戒不要贪爱这些事物：\BibleQuote{我们并不注目那看得见的，而只注目那看不见的；因为那看得见的是暂时的，那看不见的才是永远的。}\BibleRef{格后 4:18}由此可见，那些主张不但要爱，还要敬拜太阳和月亮的人，离基督徒有多远。因为若太阳和月亮不是可见的，还有什么可见的呢？但我们被禁止注目看得见的事物。因此，愿意将那份不朽的爱献给天主的人，也不可贪爱这些事物。这个问题我将在别处更详细地探讨。这里的计划不是要写关于信德，而是写关于我们藉以配得认识所信之物的生活。所以，唯独天主应当被爱；而整个世界，即一切可感觉的事物，都当被轻视——然而，它们应在此生所需时加以使用。\par

% structure: p0062 source=h2 target=section reason=relative-heading-rank; base=h2

\section{圣经谴责世俗名望和好奇求知}

38. 世俗名望在新约中这样被轻看和鄙视：圣保禄说：\BibleQuote{若我还想讨人的欢心，我就不再是基督的仆人了。}\BibleRef{迦 1:10}此外，灵魂还有一种产物，是由源自物质事物的想象而形成的，称为对事物的知识。关于这一点，我们被适当地警告要警惕好奇求知，而节制的主要功能就是纠正它。因此经上说：\BibleQuote{你们要小心，免得有人用哲学欺骗你们。}由于\Emph{哲学}一词原意是热爱和追求智慧，这是一件极具价值且应当用整个心灵去寻求的事物，宗徒非常谨慎，以免使人认为他是在阻止对智慧的热爱，所以加上了\BibleQuote{和世俗的妄言}\BibleRef{哥 2:8}。因为有些人，忽略德行，不知道天主是什么，也不知道始终如一的本性之伟大，却以为自己在以最大的好奇和热忱探究这个我们称为世界的物质团块时，是在从事一项重要的工作。这产生如此大的骄傲，以至于他们自视为他们所常谈论的天上的居民。因此，灵魂若打算为天主保持贞洁，必须戒绝这种虚妄知识的欲望。因为这种欲望通常产生错觉，使灵魂认为只有物质的东西存在；或者，如果由于权威而承认有非物质的存在，也只能以物质的形象来思想它，并且除了感官所强加的信条外，别无信仰。我们可以将关于逃避偶像崇拜的诫命应用于此。\par

39. 对于新约中要求我们不要贪爱世上任何事物的权威\BibleRef{若一 2:15}，尤其在经上说\BibleQuote{不可与此世同化}\BibleRef{罗 12:2}——因为要点是表明人爱什么，就与什么同化——对此，我若在旧约中寻找平行段落，可以找到几处；但有一卷撒罗满的书，称为训道篇，详尽地将一切属世事物视为虚空。书卷开头这样说：\BibleQuote{虚而又虚，训道者说：虚而又虚，万事皆虚。人在太阳之下辛勤劳碌，究竟有什么益处？}\BibleRef{训 1:2}如果将这些话全部思考、权衡并彻底查考，对于那些寻求逃离世界、投奔天主的人，会发现许多极其重要的东西；但这需要时间，我们的论述正急于转向其他主题。但在这开头之后，他继续详细说明，受这类事物欺骗的人是虚妄的；他称这种欺骗他们的东西为虚空——并非因为天主没有创造那些事物，而是因为人选择藉着自己的罪恶使自己屈服于那些事物，而天主的法律在善行中已使那些事物屈服于人。因为当你认为在你之下的事物是可羡慕和可欲求的，这不就是被虚幻之善所欺骗和误导吗？因此，在这些可朽而短暂的事物上节制的人，其生活准则得到两约的确认：他不可贪爱这些事物中的任何一个，也不可认为它们本身是可欲求的，而应按生活所需和职责的要求，以使用者的节制而非爱慕者的热忱来使用它们。这些关于节制的言论，相对于主题的伟大而言是简短的，但鉴于手头的任务或许又显得太多。\par

% structure: p0065 source=h2 target=section reason=relative-heading-rank; base=h2

\section{刚毅源于对天主的爱}

40. 关于刚毅，我们必须简要。我们所谈论的这种爱，应当以全部圣善之火渴望天主，它被称为节制，在于不寻求属世之物，而称为刚毅，在于承受失去它们。但在今生所拥有的一切事物中，身体，因天主的至义律法，由于古罪，是人的最沉重束缚，这一点作为事实众所周知，但其奥秘却极为难解。唯恐这束缚被动摇和扰乱，灵魂被辛劳和痛苦的恐惧所动摇；唯恐它丧失和毁灭，灵魂被死亡的恐惧所动摇。因为灵魂出于习惯而爱它，不知若能善用并明智地使用它，其复活与更新将在天主的帮助和旨意下，毫无困难地服从于它的权威。但当灵魂完全转向天主，怀着这种爱时，它便知道这些事，因此不仅会漠视死亡，甚至会渴望死亡。\par

41. 然后是与痛苦的巨大搏斗。但没有任何事物，即使是铁般的坚硬，是爱火所不能制服的。当心灵在这爱中被提升到天主时，它将以自由而光荣的姿态翱翔于一切折磨之上，带着美丽无伤的翅膀，贞洁的爱藉此升到天主的怀抱。否则，天主必定容许爱金者、爱赞美者、爱女人者比爱祂自己者更有毅力，尽管这些情况中的爱更应称为情欲或贪欲。然而，即使在这里，我们也可以看到心灵以何等力量，不顾一切惊恐，向着它所爱的对象不懈地奋力前进；我们学到，我们应忍受一切而不离弃天主，因为那些人忍受如此之多竟是为了离弃祂。\par

% structure: p0068 source=h2 target=section reason=relative-heading-rank; base=h2

\section{第二十三章——圣经中关于刚毅的诫命与范例}

42. 我在此不引用新约的权威，其中说：\BibleQuote{磨难生忍耐；忍耐生老练，老练生望德；}\BibleRef{罗 5:3} 而且除了这些话外，还有说这些话之人的事例作为证据和证实；我宁愿从旧约中引用一个忍耐的榜样，这正是摩尼教徒猛烈攻击的。我也不去提那个人，他身处极大的身体痛苦，肢体患有可怕的疾病，不仅承受了人间的苦难，还谈论了神圣的事。凡留心细听这人言论的人，将从他的每一句话学到，人们试图保持在自己权下的那些事物有何价值；而这样做时，他们自己反被情欲带入束缚，成为必朽之物的奴隶，同时又无知地寻求作它们的主人。这个人，在失去所有财富、骤然陷入极度贫穷时，他的心灵如此不动摇、专注于天主，以致表明这些事物在他眼中并不伟大，而是他相对于它们伟大，而天主相对于他伟大。\BibleRef{约 1:2} 如果在今天的人身上也能找到这种心态，新约就不会如此强烈地警告我们不要拥有这些东西以成为完善；因为拥有这些东西而不依恋它们，比根本不拥有要更值得钦佩。\par

43. 但既然我们在此谈论的是承受痛苦和身体的苦难，我就略过这个人，尽管他是如此伟大、如此不屈：这是一个男人的事例。但这些圣经向我呈现了一位具有惊人毅力的妇人，我必须立刻转到她的事例。这位妇人和她的七个孩子，宁可让暴君和刽子手从她体内取出内脏，也不让一句亵渎的话出自她口；她以劝勉鼓励她的儿子们，尽管她因他们身体所受的酷刑而痛苦，她自己也要承受她所号召他们忍受的。还有什么忍耐比这更大呢？然而，我们为何要惊讶于天主之爱，植根于她内心深处，竟能抵挡暴君、刽子手、痛苦、性别和骨肉之情呢？难道她没有听过：\Quote{在上主眼中，圣人的死是珍贵的？} 难道她没有听过：\BibleQuote{忍耐者胜过勇士？}\BibleRef{箴 16:32} 难道她没有听过：\BibleQuote{凡派定给你的，你要接受；在痛苦中要忍受；在卑微中要持守你的忍耐：因为金银在火中试炼？}\BibleRef{德 2:4} 难道她没有听过：\BibleQuote{火试炼陶工的器皿，而困苦的试炼是对义人？}\BibleRef{德 27:6} 这些她都知道，还有许多写在那些书卷中的关于刚毅的诫命，那些书是当时唯一存在的，由同一圣神所写，那圣神也书写新约中的诫命。\par

% structure: p0071 source=h2 target=section reason=relative-heading-rank; base=h2

\section{第二十四章——论正义与明智}

44. 那属于天主的正义又如何呢？正如主说：\BibleQuote{你们不能事奉两个主人。}\BibleRef{玛 6:24} 宗徒也斥责那些事奉受造物而非造物主的人，\BibleRef{罗 1:25} 旧约中不是早说过：\BibleQuote{你要朝拜上主你的天主，唯独事奉祂？}\BibleRef{申 6:13} 对此我无需再多说，因为这些书卷充满此类段落。那么，我们所描述的爱者，将从正义获得这一生活准则：他必须以完全的乐意事奉他所爱的天主，至善、至智、至和平；至于其他一切事物，他要么作为属己的来统治它们，要么以使之从属为目的来处理它们。这一生活准则，正如我们所表明的，得到两约权威的认可。\par

45. 我们也必须同样简洁地论述明智，明智的职责是分辨什么是该渴求的，什么是该避免的。没有它，我们前面所论的一切都无法实行。明智的一部分是以最警觉的戒备来守望，免得有任何邪恶的影响悄悄潜入我们心中。因此，主常大声说：\BibleQuote{醒寤。}\BibleRef{玛 24:42} 又说：\BibleQuote{你们趁着有光的时候行走，免得黑暗笼罩了你们。}\BibleRef{若 12:35} 随后又说：\BibleQuote{你们不知道少许的酵母能使整个面团发酵吗？}\BibleRef{格前 5:6} 从旧约中，没有哪段话比先知所说的\BibleQuote{轻视小事的人，必渐渐跌倒}\BibleRef{德 19:1} 更明确地谴责这种心灵的昏睡，它使我们感觉不到毁灭一步步向我们逼近。关于这个话题，倘若时间允许，我可以详加论述。假若当前的任务需要，我们或许可以证明这些奥秘的深度；亵渎者因完全无知而嘲笑它们，他们的跌倒当然不是渐进的，而是颠覆性的。\par

% structure: p0074 source=h2 target=section reason=relative-heading-rank; base=h2

\section{第二十五章——论爱天主的四项道德责任；这爱的赏报是永生与认识真理}

46. 关于正确的行为，我无须再多言。因为如果天主是人的至善——这是你无法否认的——那么，既然追求至善就是生活得好，显然生活得好无非就是全心、全灵、全意爱天主；并由此产生：这种爱必须保持完整无缺，这是节德的部分；它不因任何困难而退缩，这是勇德的部分；它不事奉任何其他对象，这是义德的部分；它警醒地审察事物，以免诡诈或欺骗潜入，这是智德的部分。这就是人的唯一成全，唯有藉此他才能成功达到真理的纯洁。两约共同命令这一点；双方同样推崇这一点。你为什么继续对你所不知的圣经加以斥责？难道你看不出你攻击这些书的愚昧吗？这些书只有那些不明白的人才会指责，而只有那些指责的人才会不明白。因为敌人不能认识它们，而认识它们的人不可能不是它们的朋友。\par

47. 那么，凡以永生为目的的人，就让我们全心、全灵、全意爱天主吧。因为永生包含了全部赏报，我们因这应许而欢喜；赏报不能先于功德，也不能在人配得之前赐予。还有什么比这更不公义，而天主又比什么更公义呢？我们不应在配得赏报之前就要求它。在此，或许问一句什么是永生并非不合时宜；不如让我们听听永生的赐予者所说的话：祂说：\BibleQuote{永生就是：使他们认识你，唯一的真天主，和你所派遣的耶稣基督。}\BibleRef{若 17:3}。所以，永生就是对真理的认识。看，那些认为传授对天主的认识就能使我们成全的人是多么乖谬和颠倒啊，因为认识真理是已经成全的人的赏报！那么，我们除了首先以完整的爱慕之情去爱那位我们渴望认识的主之外，还有什么可做的呢？由此产生我们一直坚持的原则：在天主教会中，没有什么比在论证之前运用权威更有益的了。\par

% structure: p0077 source=h2 target=section reason=relative-heading-rank; base=h2

\section{第26章——爱我们自己及爱我们的近人}

48. 继续讨论余下的事。可能会有人认为这里没有提及人自身——这位爱者。但这样想表明缺乏清晰的洞察。因为一个爱天主的人不可能不爱自己。只有那勤勉地追求至高的真福的人，才恰当地爱自己；而如果这至福除了天主之外别无他物，正如已经表明的，那么什么能阻止一个爱天主的人也爱自己呢？再者，在人群之间难道不应有互爱的纽带吗？确实如此；因此，我们想不出比人与人之间的爱更可靠的通往爱天主的步骤了。\par

49. 那么，愿主在回答关于生活诫命的问题时，赐予我们另一条诫命；因为祂不满足于一条，祂知道天主是一回事，人是另一回事，两者之间的差别正如创造者与被造为祂肖像的受造物之间的差别。祂接着说第二条诫命是：\BibleQuote{你应当爱近人如你自己。}\BibleRef{玛 22:39}。当你爱天主胜过爱你自己时，你恰当地爱自己。那么，你为自己所追求的目标，也当为你的近人追求，即：愿他以完美的爱慕之情爱天主。因为除非你试图把他引向你本人所追求的同一至善，你并不爱他如同你自己。因为这唯一的至善有空间容纳所有人与你一同追求。从这条诫命产生了人类社会中的责任，在这些责任中很难避免错误。但首先要追求的是，我们应当心怀善意，即：我们对他人不怀恶意或邪念。因为人最亲近的近邻就是人。\par

50. 也听听保禄所说的话：\BibleQuote{爱近人}，他说，\BibleQuote{不作恶。}\BibleRef{罗 13:10}。这里引用的证言很短，但若我没弄错，它们切中要点，足以说明问题。每个人都知道，在这些书中关于爱近人的话何其多、何其重要。但既然一个人可能以两种方式得罪他人：伤害他，或在有能力帮助时不去帮助他；并且由于这些事，没有一个有爱心的人会去做，因此人们被称为恶人；那么，我认为这些话语已足够证明：\BibleQuote{爱近人不作恶}，我们可以由此上升到另一句话：\BibleQuote{我们知道，一切事都协助那些爱天主的人获得益处。}\BibleRef{罗 8:28}。\par

51. 但在某种意义上，二者一同增长至圆满和成全，或者，虽然爱天主在开端是首先的，但爱近人却更先达到成全。因为起初神圣之爱或许更迅速地抓住我们，但我们在较低的事物中更容易达到成全。无论如何，要点是：没有人应当以为，他轻视近人却能得到幸福且与所爱的天主相遇。但愿，对于一个受过良好训练且心地善良的人来说，寻求近人的益处或避免伤害他，如同爱近人本身一样容易！这些事不仅需要善意，还需要高度的审慎和明智，这只有那些从万善之源的天主那里领受的人才拥有。关于这个主题——我认为这是一个极其困难的主题——我将根据我的计划尽力说几句话，将所有的希望寄托于那赐予这些恩赐的主。\par

% structure: p0082 source=h2 target=section reason=relative-heading-rank; base=h2

\section{第27章——善待我们近人的身体}

52. 人在其同胞眼中，是一个理性的灵魂，配有一个有死有肉的躯体为其服务。因此，爱近人的人，其行善部分关乎人的肉身，部分关乎其灵魂。有益于肉身者称为医药，有益于灵魂者称为纪律。此处医药包括一切维护或恢复身体健康之物。因此，它不仅包括严格意义上的医术，还包括饮食、衣服、住所，以及一切用以遮盖和保护我们身体免受内外伤害与灾祸的手段。因为饥饿、干渴、寒冷、炎热以及一切外来的暴力，都会损害我们所要考虑的健康。\par

53. 因此，那些适时而明智地提供一切所需以抵御这些邪恶和痛苦的人，被称为有怜悯心的人，尽管他们可能智慧到在行怜悯时内心不受痛苦情感的搅扰。无疑，\Quote{有怜悯心}一词意味着同情他人悲伤者心中的痛苦。同样真实的是，智者应当免于一切痛苦情感，当他帮助困乏之人、给饥饿者食物、给口渴者水喝、给赤身者衣服穿、接待旅客、释放受压迫者，最后，当他向死者施以爱德而埋葬他们时，他仍是一个恰当的有怜悯心的人，即使他心境平和地行事，并非出于痛苦情感的刺激，而是出于爱德的动机。当其中没有激情时，\Quote{有怜悯心}一词并无害处。\par

54. 相反，有些愚人，因未充分被责任感所动而不感受痛苦之情，便将行怜悯视为恶习而避免，结果冻成冷酷无情，这与理性的宁静截然不同。另一方面，天主被恰当地称为有怜悯心的；祂如何是，那些因虔诚和勤勉而堪当理解的人将会明白。恐怕我们在使用学者的言辞时，会使未受过教育的人因远离怜悯而硬起心肠，而不是以爱德之心的渴望软化他们。因此，正如怜悯要求我们为他人抵御这些痛苦，无害则禁止我们施加它们。\par

% structure: p0086 source=h2 target=section reason=relative-heading-rank; base=h2

\section{章28. 论行善于近人的灵魂。纪律的两部分：管束与训导。通过善行我们达至认识真理}

55. 至于纪律，它用以恢复心灵的健康——若无此健康，肉身健康对避免痛苦毫无助益——这是一个极其艰深的主题。正如我们在肉身上所说，治愈疾病和创伤（只有少数人能妥善处理）是一回事，而满足饥渴的需求，并在其他一切方面随时帮助他人，则是另一回事；同样，在心灵中，有些情况并不太需要教师的高超罕见职分——例如，在劝告和鼓励他人给予上述肉身所需的物品时。给予这样的劝告，就是通过纪律来帮助心灵，正如给予物品本身是通过我们的资源来帮助肉身。但还有一些情况，心灵的各种疾病以奇特而难以言表的方式被治愈。除非祂的良药从天而降，临到人类——因为他们在罪恶中如此漫不经心——否则就没有救恩的希望；事实上，就连肉身健康，追根溯源，也只能来自天主，祂赐予万物存在与福祉。\par

56. 因此，这纪律——即心灵的医药——根据我们从圣经中所能了解的，包括两件事：管束与训导。对于接受纪律者而言，管束意味着恐惧，训导意味着爱；而在施予益处者身上，有爱而无恐惧。在这两者中，天主自己——因祂的良善和仁慈，我们才得以存在——在两个盟约中赐予了我们纪律的准则。因为虽然两者在二约中都可找到，但恐惧在旧约中突出，爱在新约中突出；宗徒称前者为奴役，后者为自由。要谈论这两个盟约奇妙秩序与神圣和谐，需费时良久，许多虔诚博学之士已对此作了论述。这一主题需要多卷著作才能尽人力所能予以阐述和解释。因此，凡爱近人者，尽力在肉身上和灵魂上谋求其安全，以心灵的健康作为对待肉身的标准。至于心灵，他的努力按此顺序：首先敬畏天主，然后爱天主。这才是真正的卓越行为，由此我们获得所一直追求的对真理的认识。\par

57. 摩尼教同意我关于爱天主和爱近人的责任，但他们否认这在旧约中有教导。我认为，上面引用的关于这两项责任的经文清楚显示他们在这方面的错误有多大。但是，用一个只有极端的疯狂才能反对的简单问题：他们难道看不出这些他们所称赞的两条诫命是主从旧约中引用的，即\Quote{你当全心、全灵、全意爱上主你的天主}和另一条\Quote{你当爱人如己}？或者，如果真理之光太强，使他们不敢否认这一点，那么让他们否认这些诫命是有益的；让他们，如果可能的话，否认它们教导了最好的道德；让他们断言爱天主或爱我们的近人不是责任；让他们断言一切事不使爱天主的人得益；让他们断言爱近人不会造成伤害（这是对人类生活最有益和最卓越的双重规范）是不真实的。这样的断言不仅使他们与基督徒隔绝，而且与人类隔绝。但如果他们不敢这样说，而必须承认这些诫命的神性，那么他们为什么不停止用可怕亵渎的方式攻击和诋毁这些诫命所出自的书籍呢？\par

58. 他们是否像平常那样说，虽然我们在这些书中发现了这些诫命，但并非书中所有的内容都是好的？我不知如何应对和反驳这种狡辩。我是否要逐条讨论旧约中的话，向固执无知的人证明它们与新约完全一致？但何时才能完成？我何时有时间，他们何时有耐心？那么，该怎么办呢？我是否要放弃这事业，让他们在一个虽然虚假和邪恶却难以反驳的意见中逃脱侦察？我决不如此。天主必亲自就近援助我，他不会让我在那些困境中孤立无援或被遗弃。\par

% structure: p0091 source=h2 target=section reason=relative-heading-rank; base=h2

\section{第29章　论圣经的权威}

59. 那么，请听，你们摩尼教徒，如果你们中有些人还为你们的迷信所束缚，以致还能逃脱的话。请听，我说，不要固执，不要抗拒，否则你们的决定将致命。没有人能怀疑，你们也没有迷失到不明白：如果爱天主和爱近人是好的（这是所有人都同意的），那么依赖于这两条诫命的任何事物都不能被正确地宣称为坏。想知道依赖于它们的是什么，想从我这里学习是荒谬的。请听基督亲自说的；请听基督，我说，请听天主的智慧：他说：\Quote{这两条诫命，是全部法律和先知所依赖的。} \BibleRef{玛 22:40}\par

60. 最无耻的固执对此能说什么呢？说这些不是基督的话吗？但它们作为他的话写在福音中。说这记载是假的吗？这不是最亵渎神明的亵渎吗？这样说不是最狂妄吗？不是最鲁莽吗？不是最罪恶吗？崇拜偶像的人，甚至恨基督的名，也从未敢这样攻击这些圣经。因为，如果一件为如此强大的民众信仰所支持、为如此多的人和如此多时代的一致见证所确立的事，被置于这样的怀疑之下，以至于不被允许享有普通历史的信用和权威，那么所有文献都将彻底被推翻，所有从古代流传下来的书籍都将终结。总之，如果你引用的任何著作与我意见和目的相左，你又能引用什么呢？我也可以同样这样说。\par

61. 他们禁止我们相信一本广为流传、如今人手一册的书，却坚持要我们相信他们所引用的书，这不是无法容忍吗？如果任何著作值得怀疑，那么还有什么比一本没有名声的书，或者可能完全是伪托假名的伪书更值得怀疑呢？如果你强迫我接受这样的著作，并炫耀权威驱使我信从，那么当我拥有一本我看到长久以来广泛传播、并得到遍布全世界的教会一致见证认可的书时，我难道要因怀疑而贬低自己，更糟的是，因你的暗示而怀疑吗？即使你提出其他读本，除非有普遍同意支持，我也不会接受。既然如此，你现在除了自己愚蠢和不经思考的陈述之外，提不出任何与经文可比的东西，你认为人类会那样不合理，那样被天主眷顾所遗弃，以至于宁愿选择那些圣经——不是你所引用来驳斥的其他著作，而仅仅是你自己的话吗？你应该拿出另一份内容相同、但无讹误、更正确的手稿，只缺少你指控为伪造的段落。例如，如果你认为\WorkTitle{保禄致罗马人书}是伪造的，你就必须拿出另一份无讹误的，或者更确切地说，另一份同一宗徒的同一书信的手稿，没有错误和窜改。你说你不会，以免被怀疑窜改。这是你们惯常的答复，也是真实的。如果你这样做，我们当然会有这种怀疑；而且所有有理智的人也都会这样。那么，看看你们对自己的权威应该作何感想；并考虑是否应该相信你们的话来反对这些圣经，因为仅仅是由你拿出的手稿这一事实，就使得相信它变得危险。\par

% structure: p0095 source=h2 target=section reason=relative-heading-rank; base=h2

\section{第30章　教会作为一切智慧的导师被呼求。天主教的教义}

62. 然而何必多言？因为谁看不见，那些敢于如此攻击圣经的人，即使不是他们被怀疑的那种人，至少也不是基督徒？因为基督徒被赋予这样的生活准则：我们当全心、全灵、全意爱上主我们的天主，并爱人如己；因为全部法律和先知都系于这两条诫命。因此，天主教会，基督信徒最真实的母亲，你不仅教导唯独天主——找到祂便是最幸福的生命——必须以完美的纯洁和贞洁来钦崇，不引入任何受造物作为我们应当侍奉的崇拜对象；并且从那不朽且不可侵犯的永恒——唯独人应服从于它，唯独依附它，理性灵魂才能脱离悲惨——排除一切受造物、一切可变的、一切受时间支配的事物；而不混淆永恒、真理与和平本身所分别的，也不分离共同威严所结合的：而且你也如此包含对近人的爱与爱德，以至于灵魂因他们的罪而遭受的各种疾病，在你那里都能找到卓有成效的良药。\par

63. 你的培育和教导对孩童是稚嫩的，对青年是有力的，对老人是平安的，顾及心灵和身体的年龄。你将妇女置于丈夫之下，以贞洁和信实的服从，不是为了满足情欲，而是为了生育后代和家庭共融。你给予男人对妻子的权柄，不是为了嘲弄较弱的性别，而是出于真诚之爱的法律。你使儿女以一种自由的束缚顺服父母，并使父母以虔诚的统治管理儿女。你将兄弟以超越血缘的更强更紧密的宗教纽带联结。在不违背本性和选择的关系下，你将所有亲属关系和姻亲联盟带入彼此相爱的纽带中。你教导仆役因对工作的喜爱而非地位的必需来依恋主人。你使主人因敬畏共同的主宰天主而宽容仆役，更倾向于规劝而非强迫。你联合公民与公民，民族与民族，甚至人与人，因追念他们最初的父母，不仅在社会中，而且在兄弟情谊中。你教导君王谋求人民的福祉；你劝告人民服从他们的君王。你仔细教导谁应得尊敬，谁应得关注，谁应得敬畏，谁应得惧怕，谁应得安慰，谁应得劝诫，谁应得鼓励，谁应得管教，谁应得责备，谁应得惩罚；表明并非所有这一切都该给所有人，而是爱该给所有人，伤害不该给任何人。\par

64. 然后，在这样的人民之爱滋养并增强了依附于你怀抱的心灵，使其适于跟随天主之后，当天主的威严开始显现，正如一个在地上的居民所足够领受的那样，便产生了如此炽热的爱德，点燃了如此神圣的爱火，以致烧尽一切恶习，并净化圣化人灵，使这些话显明是神圣的：\Quote{我是吞噬的火}，以及\Quote{我来是为把火投在地上}\BibleRef{路 12:49}。这两句话来自同一天主，铭刻在旧约和新约上，以和谐的见证显示出灵魂的圣化，指向那也在新约中引自旧约的成就：\Quote{死亡被胜利吞灭了。死亡，你的刺在哪里？死亡，你的争战在哪里？}若这些异端者能理解这一句话，不再骄傲而全然和好，他们就不会在任何其他地方，而只在你的怀抱中钦崇天主。在你那里，神圣的诫命理所当然地被广泛分散的群众遵守。在你那里，理所当然地清楚明白，当法律被知晓时，罪比不知时更为严重。因为\Quote{死亡的刺就是罪，而罪的能力就是法律}\BibleRef{格前 15:56}，这增加了因明知违背诫命而击打和杀戮的力量。在你那里，理所当然地看到，在法律下的努力是多么徒劳，当情欲蹂躏心灵，因惧怕惩罚而被抑制，而非被对德行的爱所压倒。你那里，理所当然地有许多好客的、许多友善的、许多慈悲的、许多博学的、许多贞洁的、许多圣人、许多对天主的爱如此炽烈的人，以至于在完全的节制和对这世界惊人的超脱中，即使在独居中他们也找到幸福。\par

% structure: p0099 source=h2 target=section reason=relative-heading-rank; base=h2

\section{第三十一章 独修者与团修者的生活对比摩尼教徒的节制}

65. 那些能在看不见同类（尽管并非不爱他们）的情况下生活的人，他们所见的必定是何等超凡的事物？必定是某种超越人性的事物，在默观它时，人可以不见同类而生活。现在，你们这些摩尼教徒\OriginalTerm{摩尼教徒}{Manichæans}，请听完美的基督徒的习俗和显著的节制，他们不仅称赞，而且实践了贞洁的顶峰，这样，如果你们还有丝毫羞耻，就可以抑制自己在无知者面前夸耀你们的禁欲，好像它是最艰难的事一样。我要谈论一些你们并非不知、却对我们隐瞒的事情。因为谁不知道，在全世界上，尤其在东方和埃及，每日都有越来越多的完全节制的基督徒男子，正如你们不能不晓得的那样？\par

66. 我对那些方才提及的人暂不置论，他们完全避开人的视野，居住在极其荒芜的地区，满足于定期送来的简单面包和清水，与天主共融，以纯净的心灵依附于祂，并在默观祂的美善中享受至极的福乐，这美善唯有圣人的理智才能看见。我对他们暂不置论，因为有些人认为他们过度地放弃了人性的事物，却没有考虑到，那些我们无法见到其身体的人，如何能在心灵上藉祈祷、在生活中藉榜样惠益我们。但讨论这一点将耗时且无益；因为如果一个人不自动地视此种高超的圣德为可敬可佩，我们的言语又怎能令他如此呢？只是该提醒那些一无所夸的摩尼教徒：天主教会伟大圣人们的克己和节制已达到如此程度，以至有些人认为应当加以约束并召回人性的限度——他们的心神已如此高飞超乎常人，甚至在那些不赞同者眼中也是如此。\par

67. 然而，若此事超越我们的忍耐限度，那么，谁能不钦佩并称赞那些人呢？他们轻看并弃绝今世的享乐，在极其贞洁圣善的团体中共同生活，联合一致，以祈祷、阅读、讨论度日，毫无骄傲之膨胀、争辩之喧嚷或嫉妒之愠怒；反而安静、谦逊、和平，他们的生活是完美的和谐与对天主的奉献，是献于那赐予行此事能力者最悦纳的祭献。没有人拥有任何私产；没有人成为他人的负担。他们亲手劳作，从事那些既能养活身体又不使心神远离天主的职业。他们劳动的成果交给\OriginalTerm{什一监督}{decans}或\OriginalTerm{什一管理者}{tithesmen}——如此称呼乃因他们被指派管理什一之物——以便无人需为身体的照顾操心，无论在食物、衣物或任何日用所需或常见疾病方面。这些什一监督精心安排一切，迅速满足该生活方式因身体软弱而产生的需求，他们有一位称为\Quote{父亲}的人，向他呈报账目。这些父亲不仅品行更为圣善，而且通晓神圣学识，各方面德高望重；他们并无骄傲地管理那些他们称为子女的人，自己拥有发号施令的巨大权威，并得到下属欣然服从。日暮时分，他们从各自的住所聚集，在晚餐前听父亲讲道，聚集的人数至少三千之众，归于一位父亲；因为一位父亲甚至可有远超过此数的人。他们以惊人的热忱、在完全的静默中聆听，并以呻吟、眼泪或无声喜悦的表示来表达因讲道者话语而激动的心灵感受。然后，身体获得所需之补给，恰如健康与良好身体状况所需，每个人都克制不法的欲念，以免在所供应的简陋廉价食物上过度。因此，他们不仅为制伏情欲而禁绝肉食与酒，还禁绝那些更易激起口腹之欲的东西（有些人称之为\Quote{更洁净}的东西），这往往成为因追求精美食物（远非动物性食物）而不当口味的可笑可耻借口。他们拥有超出自身所需之物（由于勤劳节俭，获得甚多），便更加谨慎地分施给穷人，较之为自己获取时更为用心。因为他们不努力获取富裕，却竭力防止富余留于自己手中——甚至运载整船货物至贫民居住之地。此事尽人皆知，我无需赘言。\par

68. 妇女的生活也是如此，她们勤勉贞洁地事奉天主，按照合宜的要求与男子分开而居，只在虔诚之爱与德行的模仿上与男子联合。不允许青年男子接近她们，甚至年长者，无论多么可敬与认可，也不得进入，除非在门廊处，以便提供必要物品。因为妇女以纺羊毛维持生计，并将布匹交给弟兄们，从他们那里换取所需食物。这样的习俗、这样的生活、这样的安排、这样的体系，即使我想要称赞，也不能恰如其分地称赞；此外，我担心，倘若我认为有必要在单纯的叙述之上添加华丽的颂词，似乎会显得我以为单凭描述不足以使人接纳。你们摩尼教徒，若能，则在此处吹毛求疵吧。不要将我们的莠子显扬在那些过于盲目无法分辨的人面前。\par

% structure: p0104 source=h2 target=section reason=relative-heading-rank; base=h2

\section{第32章 赞美神职人员}

69. 然而，天主教会道德上的卓越并非如此狭隘，以致我应将其赞美仅限于此处提及的那些人的生活。因为我认识多少极其卓越圣善的主教，多少司铎，多少执事，以及各种神圣圣事的服务者，他们的德行在我看来更令人钦佩、更值得称赞，因为在纷繁多样的人群中、在这动荡的生活中保持德行更为困难！因为他们所管理的人既有需要治愈的，也有已经治愈的。群众的恶习必须忍受，以便他们得以治愈，瘟疫必须在被制服之前予以忍受。在此保持最优的生活方式和平静安宁的心灵极为艰难。简言之，我们在此是与正在学习生活的人群相处；而那里的人则是在生活。\par

% structure: p0106 source=h2 target=section reason=relative-heading-rank; base=h2

\section{第33章 另一类共同生活在城中的人。三天的斋戒}

70. 然而，我不会因此轻视一群值得称赞的基督徒——即那些共同生活在城市中、完全脱离世俗生活的人。我在米兰见过一所圣人的居所，人数不少，由一位司铎主持，他极有德行和学识。在罗马，我认识几个地方，每处都有一位以品德庄重、审慎和神圣知识著称者，主持一切与他共同生活的人，彼此生活在基督徒的爱德、圣洁和自由中。这些人也不成为任何人的负担；而是按东方的习俗，并依宗徒保禄的权威，亲手劳动维持生计。我听说许多人进行令人惊异的严厉斋戒，不仅每天只吃一顿晚餐（这在各处相当普遍），而且常常连续三天或更长时间不进饮食。这不仅在男人中，也在妇女中，她们大量以寡妇或贞女身份共同生活，靠纺纱织布为生，在每一处都有一位最有判断力和经验、不仅善于指导及塑造道德行为，而且善于教导理智的妇人主持。\par

71. 凡此种种，没有人被迫承受自己能力不及的艰苦；没有任何事违反人意强加于人；也没有人因承认自己太软弱无法效法他们而受其余人责难：因为他们牢记圣经如何强烈命令众人爱德；他们牢记\BibleQuote{为洁净者一切都洁净}\BibleRef{铎 1:15}，以及\BibleQuote{不是入口的使人污秽，而是出口的}\BibleRef{玛 15:11}。因此，他们一切努力关心的不是拒绝某种食物为不洁，而是克制不当的欲望并保持弟兄友爱。他们记得：\BibleQuote{食物是为肚腹，肚腹是为食物；但天主将两者都要消灭}\BibleRef{格前 6:13}；又说：\BibleQuote{我们不吃，也不欠缺；我们吃，也不多余}\BibleRef{格前 8:8}；而最重要的是：\BibleQuote{我的弟兄们，不吃肉，不饮酒，不作任何使弟兄跌倒的事，是好的}；因为这段经文表明所有这一切都以爱德为目的。他说：\BibleQuote{有人相信可吃一切；软弱的人只吃蔬菜。吃的人不可轻视不吃的人；不吃的人也不可判断吃的人，因为天主已接纳了他。你是谁，你竟判断别人的仆人？他或站或倒，只对自己的主人；他必站立，因为天主能使他站立。}稍后又说：\BibleQuote{吃的人为主吃，并感谢天主；不吃的人为主不吃，也感谢天主。}在后面又说：\BibleQuote{所以我们每人都要将自己的事向天主交代。因此我们不要再彼此判断；宁可判断这点：不给弟兄放置绊脚石或使人跌倒的机会。我在主耶稣内知道，并深信没有一样东西本身是不洁的；但若有人以为某物是不洁的，那对他就是不洁的。}他岂不是更清楚地表明，不在于我们吃的东西，而在于意念，才有能力使之污秽；因此，即使那些能轻视这些东西、完全知道自己在精神优越中取用任何食物而不成为贪食者时并不会被污秽的人，也仍需顾念爱德？请看他所加的：\BibleQuote{如果你的弟兄因你的食物而忧愁，你便不是按爱德行事。}\BibleRef{罗 14:2}\par

72. 其余的部分请自行阅读：引用全部太长了。你会看到，那些能轻视这类事物的人——即那些力量较大、根基稳固的人——被告知仍须戒避，以免那些因软弱仍需此种戒避的人跌倒。我所描述的人知道并遵守这些事；因为他们是基督徒，而非异端者。他们按宗徒的教导理解圣经，而非依据自命不凡、虚构的宗徒之名。不吃的人没有人轻视；吃的人没有人判断；软弱的人吃蔬菜。然而许多强壮的人为了软弱者的缘故也这样做；许多人这样做的原因不是这个，而是为了廉价的伙食，并以最低的生活费用过最平静的生活。他说：\BibleQuote{凡事我都可行，但我不要受任何事物的辖制。}\BibleRef{格前 6:12}因此许多人禁食肉类，却不迷信地视之为不洁。许多人不饮酒，但他们不以为酒污秽他们；因为他们以极大的得体与节制将酒给予虚弱的人——总之，给予所有没有酒便不能保持身体健康的人。当有人愚昧地拒绝酒时，他们以弟兄的身份劝告他们，不要让愚蠢的迷信使他们更虚弱，而不是更圣洁。他们为他们读宗徒对弟子的训诫：\BibleQuote{为了你的许多病痛，稍微饮点酒。}\BibleRef{弟前 5:23}然后他们勤修虔诚；他们知道\BibleQuote{身体的锻炼略有助益}\BibleRef{弟前 4:8}，如同同一宗徒所说。\par

73. 因此，那些有能力的人（他们不计其数）戒食肉类与酒，出于两种原因：或为了弟兄的软弱，或为了自己的自由。他们主要关注爱德。在饮食上有爱德，在言语上有爱德，在衣着上有爱德，在目光上有爱德。爱德是他们相遇的焦点和行事的准则。违背爱德被认为如同违背天主一样有罪。凡反对爱德的，便受到攻击并逐出；凡伤害爱德的，不容许存在一日。他们知道这是基督和宗徒们所命令的；没有爱德，一切都是空的；有了爱德，一切都成全。\par

% structure: p0111 source=h2 target=section reason=relative-heading-rank; base=h2

\section{第三十四章．不应因坏基督徒、坟墓与画像崇拜者的行为而责备教会}

74. 你们这些摩尼教徒，如果你们能，就对这些提出异议吧。看看这些人，并若无谎言地指责他们。比较你们的斋戒与他们的斋戒，你们的贞洁与他们的贞洁；在衣着、饮食、自制，最后，在爱德上，将你们与他们相比。比较他们的诫命与你们的诫命，这是最切要的。那么你们就会看到做作与真诚、正路与邪路、信德与欺骗、力量与浮夸、福乐与悲惨、合一与分裂之间的区别；简言之，就是迷信的塞壬与宗教的港口之间的区别。\par

75. 不要向我招来那些基督徒名称的教授者，他们既不知道也不证实他们职业的能力。不要搜罗许多无知之人，他们在真宗教中也迷信，或者如此沉溺于邪恶情欲以致忘记他们对天主许下的诺言。我知道有许多坟墓和画像的敬拜者。我知道有许多人在死者面前纵饮至醉，并在他们为尸体举办的宴会中，将自己埋葬于被埋葬者之中，并以宗教之名给自己的贪食和酗酒。我知道有许多人口头上弃绝了这个世界，却仍愿被世俗事物的一切重量所累，并以此为乐。在如此众多的人群中，你们能找到一些人，通过谴责他们的生活来欺骗粗心之人并引诱他们脱离天主教的安全，这并不奇怪；因为在你们少数人中，当被要求展示哪怕一个你们所谓的\OriginalTerm{特选者}{the elect}遵守诫命时，你们却不知所措，而你们在自己不可辩护的迷信中却宣称这些诫命。这些人是多么愚蠢、多么不虔敬、多么有害，以及他们被你们中大多数人——几乎所有人——如何忽视，我在另一卷书中已表明。\par

76. 我现在对你们的劝告是：你们至少应停止诽谤天主教会，不要攻击教会在谴责的那些人的行为；教会如同对待邪恶的子女一样，每日寻求纠正他们。然后，若他们中有人凭着善意并靠着天主的帮助得到改正，他们便以补赎重新获得因罪而失去的东西。而那些怀着恶意坚持旧恶习，甚至添加更坏恶习的人，固然被允许留在主的田地里，与好种子一同生长；但分离莠子的时刻终将到来。或者，若由于他们至少拥有基督徒的名称，他们应被置于糠秕而非蓟草之中，也会有那一位来扬净禾场，将糠秕与麦子分开，并按照各自的功过给予应得的赏报。\par

% structure: p0115 source=h2 target=section reason=relative-heading-rank; base=h2

\section{第三十五章——宗徒准许已受洗者享有婚姻与财产}

77. 同时，你们为何发怒？为何党派精神蒙蔽你们的眼睛？为何你们陷入如此巨大错误的冗长辩护中？在田地中寻找果实，在禾场上寻找麦子：它们会轻易被找到，并向寻求者呈现。为什么你们只专注于渣滓？为什么你们利用篱笆的粗糙来吓唬无经验者，使其远离园子的肥沃？有一个恰当的入口，虽然只有少数人知道；通过它人们进入，虽然你们不信或不愿寻找。在天主教会中，有无数信徒不使用这个世界，也有那些人\Quote{使用它}，按宗徒的话说，\BibleQuote{如同不使用它}，\BibleRef{格前 7:31}，正如在基督徒被迫崇拜偶像的那些时代所证明的。那时，多少富人、多少农夫家主、多少商人、多少军人、多少自己城中的领袖、多少元老、男女众人，放弃所有这些空虚和短暂的事物，虽然当他们使用这些时并不被它们束缚，为得救的信德和宗教忍受死亡，并向不信者证明他们不是被这些事物所占有，而是真正占有它们。\par

78. 你们为何指责我们，说在圣洗圣事中更新的人不应再生育儿女，不应拥有田地、房屋和钱财？\Person{保禄}允许这样。因为，无可否认，他在列举了许多作恶者不得承受天主的国之后，写信给信徒说：\Quote{你们从前也如此，他说，但你们已洗净了，已圣化了，已因主\Person{耶稣基督}的名和藉我们天主的\OriginalTerm{圣神}{Spirit}而成义了。} 确实，无人会冒然以为那被洗净、被圣化的不是信徒，不是那些弃绝了这个世界的人。但在指明他写信给谁之后，让我们看看他是否允许他们拥有这些。他接着说：\Quote{凡事我都可行，但不都有益处；凡事我都可行，但我却不受任何事物的奴役。食物是为肚腹，肚腹是为食物；但天主必要使两者同归于尽。身体不是为淫乱，而是为主，主也是为身体。然而天主使主复活了，他也要藉自己的能力使我们复活。难道你们不知道你们的身体是基督的肢体吗？我该将基督的肢体作为娼妓的肢体吗？断乎不可！难道你们不知道，凡与娼妓结合的，便是与她成为一体吗？因为他说过：\InnerQuote{二人成为一体。}但那与主结合的，便是与主成为一神。你们要远避淫乱。人所犯的任何罪都在身体以外，但那犯淫乱的人，却是得罪自己的身体。难道你们不知道你们的身体是住在你们内的圣神的宫殿吗？这圣神是天主所赐的，你们已不属于自己；因为你们是重价买来的：所以要在你们的身体上光荣天主，并携带他。} \BibleRef{格前 6:11} \Quote{但关于你们所写的事，我认为男人不亲近女人是好的。但为了避免淫乱，男人当各有自己的妻子，女人也当各有自己的丈夫。丈夫对妻子当尽应尽的责任，妻子对丈夫也当如此。妻子对自己的身体没有主权，丈夫有；同样，丈夫对自己的身体没有主权，妻子有。你们不可彼此亏负，除非两厢情愿，暂时分房，为专心祈祷；以后仍要同房，免得撒殚因你们不能节制而诱惑你们。我说这话是出于宽容，并非出于命令。因为我愿意众人像我一样；但每人都有自己从天主得来的恩赐，有人这样，有人那样。} \BibleRef{格前 7:1}\par

79. 宗徒，你以为，是否已充分向强壮者指明了最高的事，并向软弱者允许了次好的事？不亲近女人，他指明是最高的事，当他说：\Quote{我愿众人像我一样。}但仅次于这最高的，是夫妻间的贞洁，以免人成为淫乱的猎物。他是否说过，这些人因为结了婚就不是信徒？实际上，正是藉着这夫妻间的贞洁，他说，若一方是不信者，那结合的二人彼此圣化，他们的子女也圣化了。他说：\Quote{不信主的丈夫因信主的妻子而圣化，不信主的妻子因信主的丈夫而圣化；不然，你们的子女就不洁净，但如今他们是圣洁的。} \BibleRef{格前 7:14} 你们为何要反对如此明显的真理？为何试图用虚妄的阴影遮蔽圣经的光明？\par

80. 不要说，慕道者可以娶妻，信徒却不可以；慕道者可以有钱，信徒却不可以。因为有许多人像是使用，却好像不使用。在那神圣的洗濯中，新人的更新已然开始，逐渐达致成全，有的人更快，有的人更慢。然而，若我们不怀敌意而留心观察，在许多人身上确实能见到向着新生命的进步。正如宗徒论及自己所说：\Quote{我们外在的人日渐朽坏，内在的人却日日更新。} \BibleRef{格后 4:16} 宗徒说内在的人日日更新，为能达致成全；你们却希望它一开头就达致成全！如果你们真的这样希望，倒也罢了。实际上，你们的目的不是扶助软弱的人，而是误导疏忽的人。你们不该如此傲慢地说话，即使你们在你们幼稚的训诫上被公认为完美。但当你们的良心知道，你们领入你们教派的人，在更亲密地认识你们之后，会发现你们身上有许多事，是任何听你们指责别人的人所不会怀疑的，那么，你们要求软弱的公教徒达致成全，以转移经验不足的人离开天主教会，而你们自己向那些被转移的人却丝毫不显露出这等成全，岂不是极大的无礼？但为了不显得无故抨击你们，我现在要结束这卷书，最终将开始陈述你们生活的规诫和你们著名的习俗。\par

\SourceRecord{Translated by Richard Stothert.；Nicene and Post-Nicene Fathers, First Series；Vol. 4.；Edited by Philip Schaff.；Buffalo, NY: Christian Literature Publishing Co.,；1887.；<http://www.newadvent.org/fathers/1401.htm>.}
